\chapter{Graça, fé e perseverança}
\label{cap:graca_fe}
\noindent\textit{Vinheta.} Um jovem recém-batizado decide “limpar o feed” e reduzir o tempo de tela. Na primeira semana, vai bem. Na segunda, desliza. Sente culpa, então tenta “compensar” com regras rígidas e promessas que não cumpre. Ao final do mês, conclui: “\textit{Eu não dou conta}”. 
O problema não é a \emph{falta de ferramenta}, é a \textbf{fonte} de onde ele opera. O evangelho oferece essa fonte: \textbf{graça} que recebemos pela \textbf{fé}, gerando \textbf{perseverança}.

\section{O motor da fidelidade}
O Novo Testamento é claro: “\emph{Pela graça sois salvos, mediante a fé; e isto não vem de vós, é dom de Deus}” (Ef~2:8--9). A vida cristã \textit{começa} e \textit{continua} assim: graça recebida, fé que confia, amor que trabalha (Gl~5:6). Em tempos de pressão tecnológica (vigilância, persuasão, coerção), quem persevera não é o “mais disciplinado”, mas quem \textbf{permanece em Cristo} (Jo~15).

\section{Graça que educa, não que desculpa}
Graça não é tolerância ao pecado; é o poder de Deus que nos \textbf{educa} (Tt~2:11--12). Ela nos treina a renunciar “impiedade e paixões mundanas” e a viver “sensata, justa e piedosamente” \textit{neste século}. Isso inclui nosso trato com telas, algoritmos e conveniências.

\section{Fé que confia em meio ao ruído}
Fé bíblica não é pensamento positivo, é \textbf{confiança relacional} em Deus revelado em Cristo. Em um ambiente onde muitas “verdades” disputam sua atenção, a fé pratica duas posturas: \textit{ouvir} (Palavra) e \textit{obedecer} (passos concretos). A Palavra ancora a narrativa; a obediência realinha a identidade.

\section{Perseverança: a arte de continuar}
Perseverança é permanecer sob o peso (\emph{hypomonē}). Não é teimosia cega; é \textbf{esperança aplicada} (Rm~5:3--5; Tg~1:2--4). Em termos práticos, a perseverança cristã se sustenta em três meios de graça ordinários: \textit{Palavra, oração e comunhão}. Tecnologia pode ajudar (lembretes, leituras, calendário), mas não substitui \textbf{pessoas} e \textbf{liturgias} simples.

\section{NIC aplicado: por que graça, fé e perseverança importam no digital}
\begin{nicbox}
\textbf{Narrativa} --- A graça reescreve a história: não somos \emph{o que entregamos}, somos \emph{os que recebemos}.\\
\textbf{Identidade} --- A fé nos ancora em Cristo: não somos perfis; somos filhos adotivos.\\
\textbf{Controle} --- A perseverança responde à coerção: quando portas fecham, respondemos com fidelidade criativa.
\end{nicbox}

\section{Práticas que funcionam em 15 minutos por dia}
\begin{checklist}
\begin{itemize}[leftmargin=*,noitemsep]
  \item \textbf{Lectio breve} (5 min): leia um salmo em voz alta; sublinhe um verbo de Deus.
  \item \textbf{Oração dirigida} (5 min): agradeça por 3 graças recebidas no dia; peça 1 ajuda concreta.
  \item \textbf{Passo pequeno} (5 min): um ato de obediência mensurável (ligar para alguém, pedir perdão, agendar encontro).
\end{itemize}
\end{checklist}

\section{Higiene de telas que nasce da graça}
Não começamos pelo “não pode”; começamos pelo “\textit{a quem pertencemos}”. A partir disso, implementamos limites \emph{simples}:
\begin{checklist}[title={Regra de bolso (7 dias)}]
\begin{itemize}[leftmargin=*,noitemsep]
  \item Um \textbf{primeiro olhar} do dia para a Palavra, não para o celular.
  \item Uma \textbf{janela de tela} definida (p.\,ex., 8h–21h) e o celular dorme fora do quarto.
  \item Uma \textbf{refeição sem tela} por dia com conversa e oração curta.
  \item Um \textbf{dia de sabá digital} (6h sem redes sociais).
\end{itemize}
\end{checklist}

\section{Perguntas para grupos pequenos}
\begin{itemize}[leftmargin=*,noitemsep]
  \item Onde você tem tentado “vencer” no braço aquilo que só a graça pode operar?
  \item Que prática de 15 minutos por dia você pode iniciar esta semana?
  \item Como sua igreja pode facilitar perseverança (pequenos grupos, leitura comunitária, mentorias)?
\end{itemize}

\bigskip
\noindent\textit{Próximo capítulo: Verdade em tempos de deepfake --- discernindo narrativas e testando espíritos.}
